% !TeX root = ../presentacion.tex
% ===========================================================================
%  10. Calidad y despliegue — 0:35
%
%  Criterio: «Calidad técnica y documentación» (10 %). Un atributo de calidad
%  sin mecanismo es una intención: por eso la tabla tiene dos columnas.
%
%  Cifras verificadas: docker-compose.yml define 11 servicios; el backend
%  tiene 73 métodos @Test repartidos en 12 clases de prueba.
% ===========================================================================

\bloque{Calidad y despliegue}{0:35}{\exponeArquitectura}%
       {Calidad técnica y documentación (10\,\%)}

\begin{frame}{Un comando, once contenedores, ningún requisito previo}
  \note{La cifra grande es el argumento: no hay guía de instalación de
  cuatro páginas porque no hace falta instalar nada. Que la evaluación no
  dependa de qué tenga instalado quien evalúa no es comodidad, es una
  decisión.}

  \begin{columns}[c]
    \begin{column}{0.24\textwidth}\cifra{1}{\id{docker compose up}}\end{column}
    \begin{column}{0.24\textwidth}\cifra{11}{contenedores}\end{column}
    \begin{column}{0.24\textwidth}\cifra{73}{pruebas automatizadas}\end{column}
    \begin{column}{0.24\textwidth}\cifra{0}{cosas que instalar}\end{column}
  \end{columns}

  \vspace{8pt}
  \begin{center}
    \scriptsize
    \begin{tabular}{@{}lL{7.4cm}@{}}
      \toprule
      \textbf{Atributo} & \textbf{Mecanismo concreto}\\
      \midrule
      Mantenibilidad & Dominio aislado de HTTP y de JPA; cada regla con su prueba\\
      Escalabilidad  & Servicios sin estado tras el gateway; se replican por separado\\
      Seguridad      & Bases aisladas por credenciales; el navegador nunca alcanza un servicio\\
      Disponibilidad & \id{healthcheck} por contenedor; un remote caído degrada una ruta, no la aplicación\\
      \bottomrule
    \end{tabular}
  \end{center}

  \vspace{-2pt}
  {\tiny\color{textoSuave} Ni Java, ni Node, ni Maven, ni PostgreSQL en la
  máquina: todo se compila dentro de los contenedores. Manual de despliegue
  aparte (\ent{5}).}
\end{frame}
